% =====================================================================
% SECCION 5 -- ESPECIFICACION DE REQUERIMIENTOS
% Responsable del apartado 5.2: Fabian Moreno Ugarte
%
% Creada el 14/09/2026 al reestructurar el informe al indice del modelo
% del curso. 5.1 es esqueleto. 5.2 recoge los requisitos de privacidad
% M-01 a M-10, que antes estaban en el marco normativo (hoy Anexo B, B.5),
% y el requisito de no identificacion (RNF-01) descrito en palabras. Nada
% de 5.2 afirma algo que no conste en el Anexo B: si se corrige un
% requisito, se corrige en los dos sitios.
%
% El \label{sec:privacidad-diseno} vive aqui: lo usan el Anexo A (medidas
% M-02, M-03 y M-04) y ya no lo genera md2tex.py.
%
% 24/09/2026: 5.1 deja de ser esqueleto y se anade 5.3.
%   CU-01 a CU-07, RF-01 a RF-07 y RNF-02 a RNF-06 se trasladan desde la
%   Seccion 5 de INFORME_COMPLETO_ANGEL.md, sin sus referencias internas.
%   CU-08 a CU-12 y RF-08 a RF-12 son nuevos.
%   Todos los casos llevan los cuatro campos (actor, precondicion, flujo,
%   postcondicion). Los campos que el documento de origen no traia se
%   completaron a partir del codigo del prototipo, no de la documentacion,
%   y estan pendientes de revision por Liderazgo e Integracion.
%   RNF-02 a RNF-06 van en 5.3 y no en 5.2 porque dos de ellos (RNF-05 y
%   RNF-06) no son de privacidad y porque 5.2 no debe afirmar nada que no
%   conste en el Anexo B. El RNF-01 sigue descrito en 5.2 con su marcador:
%   el documento de origen tampoco trae su texto literal.
%
% \casodeuso{ID}{titulo}{actor}{precondicion}{flujo}{postcondicion}
\newcommand{\casodeuso}[6]{%
  \subsubsection*{#1. #2}
  \textit{Actor:} #3

  \textit{Precondición:} #4

  \textit{Flujo:} #5

  \textit{Postcondición:} #6\par}
% =====================================================================
\section{Especificación de Requerimientos}
\label{sec:requerimientos}

\subsection{Requerimientos funcionales}
\label{sec:req-funcionales}

Este apartado describe el prototipo tal como está implementado, no una
versión aspiracional. Los casos de uso recogen lo que hace hoy la consola
AeroTrack y el servidor local que la atiende; la tabla final resume los
requisitos funcionales que se derivan de ellos y su estado en el
prototipo.

\casodeuso{CU-01}{Configurar el plano y las cámaras}
  {Operador.}
  {el servidor local está en ejecución y el usuario ha iniciado sesión con
   el rol de operador.}
  {el operador elige uno de los planos del aeropuerto, crea un plano en
   blanco o importa una imagen, un PDF o un archivo DXF. Si importa una
   imagen o un PDF, recorta la región útil y el sistema la reduce a un
   dibujo esquemático de líneas por detección de bordes, en lugar de
   conservar la fotografía completa. Después coloca cada cámara sobre el
   plano, define su cobertura orientativa (cono, rectángulo o forma libre)
   y declara qué cámaras son vecinas entre sí.}
  {el plano queda disponible para calibración.}

\casodeuso{CU-02}{Calibrar una cámara}
  {Operador.}
  {la cámara tiene una fuente de video válida.}
  {el operador marca un punto de referencia del suelo en el video y su
   correspondencia en el plano, un mínimo de cuatro veces; se recomiendan
   de seis a ocho puntos bien distribuidos. El sistema calcula la
   homografía y rechaza la calibración si hay referencias repetidas o casi
   alineadas. Con cinco o más puntos comprueba además cada referencia
   dejándola fuera del ajuste.}
  {la cámara queda habilitada para proyectar posiciones al plano
   compartido.}

\casodeuso{CU-03}{Procesar una fuente y hacer seguimiento}
  {Operador.}
  {la cámara tiene una fuente configurada; si se usa YOLO, el operador
   dispone de los pesos del modelo en el equipo.}
  {el operador elige un detector (HOG de referencia, P2PNet o YOLO) e
   inicia la sesión. El sistema detecta personas en cada cuadro y asigna a
   cada una un identificador local estable dentro de esa cámara. Con HOG y
   P2PNet el seguimiento es una adaptación de
   ByteTrack~\cite{Zhang_2022_ECCV} a puntos, con asignación húngara en dos
   etapas; con YOLO se sigue la caja completa de cada persona. Con P2PNet
   la posición no se proyecta al plano, porque el detector localiza
   cabezas y no pies y esta integración no aproxima una a otra sin validar
   el error que eso introduciría; ese modo sirve para validar detección y
   seguimiento por cámara, no el plano completo.}
  {cada persona detectada tiene un identificador local en esa cámara y, si
   la cámara está calibrada y el detector no es P2PNet, una posición en el
   plano.}

\casodeuso{CU-04}{Asociar identidad entre cámaras}
  {el sistema, a partir de la declaración de sincronización que hace el
   operador.}
  {al menos dos cámaras calibradas y vecinas, y el operador declaró
   expresamente que verificó la sincronización de reloj entre las fuentes.}
  {ante una observación nueva, el sistema busca entre las personas activas
   de una cámara vecina alguna cuya posición proyectada, ajustada por
   velocidad si no hay solape temporal, quede dentro de un margen de
   distancia y tiempo configurables. Si hay un candidato inequívoco,
   reutiliza su identificador y lo marca como estimado; si hay varios
   candidatos próximos, desempata con un descriptor de apariencia no
   biométrico; si ninguno es suficientemente inequívoco, lo marca como
   incierto en vez de forzar una fusión. Si la sincronización no está
   declarada, el sistema conserva los identificadores locales y no asocia
   personas entre cámaras.}
  {cada observación queda con un identificador global y un estado de
   asociación: local, estimada o incierta.}

\casodeuso{CU-05}{Detectar y alertar aglomeraciones}
  {Operador o administrador.}
  {al menos una cámara configurada y una sesión de monitoreo en curso.}
  {el usuario define, por zona o de forma global, un radio, una cantidad
   mínima de personas y un tiempo de permanencia. El sistema marca una
   alerta cuando una concentración cumple los tres criterios de forma
   sostenida.}
  {la alerta queda registrada como incidente en la vista de alertas, con
   un estado que el personal puede revisar, y se envía por correo si el
   aviso por correo está configurado.}

\casodeuso{CU-06}{Visualizar el plano en vivo y el video por cámara}
  {Operador o administrador.}
  {una sesión de monitoreo en curso.}
  {el usuario ve el plano con nodos, trayectorias recientes y
   aglomeraciones resaltadas, y puede pasar al video real de una cámara con
   la caja y el identificador de cada persona superpuestos.}
  {la consulta no modifica la configuración ni la sesión.}

\casodeuso{CU-07}{Exportar reportes}
  {Operador o administrador.}
  {existe una sesión de monitoreo.}
  {el usuario exporta el reporte de la sesión en PDF, Excel, CSV o JSON.
   El reporte indica siempre si los datos proceden de una fuente real o de
   la demostración sintética.}
  {el usuario obtiene el archivo del reporte; la sesión no se modifica.}

\casodeuso{CU-08}{Gestionar proyectos independientes}
  {Operador; el administrador solo puede abrir proyectos existentes.}
  {el usuario ha iniciado sesión.}
  {el operador crea un proyecto dándole un nombre, lo abre, lo renombra o
   lo elimina. Al entrar en la consola con más de un proyecto guardado, el
   usuario elige en cuál trabajar. El sistema rechaza los cambios
   destinados a un proyecto distinto del que está abierto, para que el
   borrador de un proyecto no sobrescriba el plano de otro.}
  {cada proyecto conserva por separado su plano, sus cámaras, sus zonas y
   su calibración.}

\casodeuso{CU-09}{Revisar la cámara en vivo antes de calibrar}
  {Operador.}
  {la cámara tiene una fuente configurada y no hay una sesión de
   seguimiento en curso. No se requiere calibración.}
  {en el paso de prueba del asistente, el operador abre la vista en vivo
   de la cámara para comprobar encuadre, luz y retardo. Si la fuente es una
   grabación, puede pausarla, reanudarla y saltar a un instante; si es una
   cámara en vivo, ve la imagen tal como llega. También puede obtener una
   imagen fija para marcar la zona útil.}
  {la apertura de la vista queda en el registro de auditoría. Al obtener
   una imagen fija, la fuente queda marcada como comprobada, con su
   resolución y su frecuencia de cuadros.}

\casodeuso{CU-10}{Configurar y recibir avisos de equipaje sin custodia}
  {Operador, que configura el aviso; operador o administrador, que lo
   recibe.}
  {la cámara tiene una fuente configurada y el equipo dispone de los pesos
   del modelo YOLO.}
  {el operador activa el aviso en cada cámara y fija el tiempo que un
   bulto debe permanecer quieto antes de avisar, entre 10 y 7200 segundos
   (90 por defecto), y cada cuánto se revisa, entre 2 y 60 segundos (4 por
   defecto). Durante el monitoreo el sistema busca mochilas, bolsos y
   maletas cada pocos segundos y, si uno permanece quieto más del tiempo
   fijado, genera un aviso de equipaje sin custodia.}
  {el aviso queda registrado como incidente en la vista de alertas y se
   envía por correo si el aviso por correo está configurado. Es un aviso
   para revisión humana, no una conclusión de seguridad: también se activa
   con el equipaje que un pasajero deja a su lado mientras espera.}

\casodeuso{CU-11}{Configurar y recibir alertas de seguridad por correo}
  {Operador o administrador.}
  {el usuario ha iniciado sesión y dispone de una cuenta de correo
   remitente con contraseña de aplicación.}
  {desde la vista de alertas, el usuario activa el aviso por correo e
   indica la cuenta remitente, su contraseña de aplicación y hasta veinte
   destinatarios, y puede enviar un correo de prueba. Ante cada incidente
   nuevo de aglomeración o de equipaje, el sistema envía un correo con el
   lugar, la cámara, el motivo del aviso y la advertencia de que requiere
   comprobación humana, y no reenvía un aviso ya enviado.}
  {la configuración queda guardada fuera de los archivos del proyecto, la
   contraseña no se muestra nunca en la interfaz y el cambio queda en el
   registro de auditoría.}

\casodeuso{CU-12}{Consultar el panel comercial}
  {Operador o administrador.}
  {existe una sesión de monitoreo.}
  {en la vista de reportes, pestaña de monitoreo integrado, el usuario
   consulta el panel comercial: zonas ordenadas por oportunidad comercial,
   contexto de la medición, traza de ocupación, tráfico por acceso y avisos
   para LAP. Puede descargar el mismo contenido en PDF, Excel, CSV o JSON.}
  {la consulta no modifica la sesión; el panel indica si los datos
   proceden de una fuente real o son sintéticos.}

La tabla siguiente resume los requisitos funcionales que se derivan de
estos casos de uso y su estado verificado en el prototipo.

{\small
\setlength{\LTleft}{0pt}\setlength{\LTright}{0pt}
\begin{longtable}{@{}L{1.4cm}L{\dimexpr\linewidth-5.4cm-4\tabcolsep\relax}L{4.0cm}@{}}
\toprule
\textbf{ID} & \textbf{Requisito} & \textbf{Estado} \\
\midrule
\endfirsthead
\toprule
\textbf{ID} & \textbf{Requisito} & \textbf{Estado} \\
\midrule
\endhead
\bottomrule
\endlastfoot
RF-01 & Configurar plano y cámaras, con recorte e importación como dibujo de
líneas & Implementado \\
\addlinespace
RF-02 & Calibrar por homografía con validación de consistencia geométrica &
Implementado \\
\addlinespace
RF-03 & Detectar y seguir personas con detector seleccionable (HOG,
P2PNet o YOLO) & Implementado; YOLO requiere pesos locales no incluidos \\
\addlinespace
RF-04 & Asociar identidad entre cámaras solo con sincronización declarada,
con estado incierto cuando corresponda & Implementado \\
\addlinespace
RF-05 & Reglas de aglomeración configurables por radio, cantidad y
permanencia & Implementado \\
\addlinespace
RF-06 & Visualizar el plano en vivo y el video real por cámara con el
identificador superpuesto & Implementado \\
\addlinespace
RF-07 & Exportar reportes en PDF, Excel, CSV o JSON & Implementado \\
\addlinespace
RF-08 & Gestionar proyectos independientes, cada uno con su propio plano,
cámaras y calibración & Implementado \\
\addlinespace
RF-09 & Revisar la cámara en vivo antes de calibrarla & Implementado \\
\addlinespace
RF-10 & Configurar por cámara y emitir avisos de equipaje sin custodia &
Implementado \\
\addlinespace
RF-11 & Configurar y enviar alertas de seguridad por correo &
Implementado \\
\addlinespace
RF-12 & Consultar el panel comercial y descargar su contenido &
Implementado \\
\end{longtable}
}

\subsection{Requerimientos de privacidad}
\label{sec:privacidad-diseno}

La normativa de protección de datos exige que las garantías se incorporen
desde la concepción del tratamiento y que la configuración más protectora
venga activada por defecto (Sección~\ref{sec:marco}). El sistema recoge
esa exigencia como requisitos no funcionales, es decir, como condiciones
que debe cumplir con independencia de su funcionalidad.

El requisito rector es que de la imagen no se deriven rasgos que permitan
identificar a una persona, ni siquiera de forma indirecta (identificado en
el proyecto como RNF-01). Su texto literal y el documento en que consta no
figuran en la documentación disponible al cerrar esta versión, de modo que
se recoge según lo reportado por el equipo y queda sujeto a contrastarlo
con su fuente.\verificar{texto literal del RNF-01 y documento en que
consta. No figura en la documentación del proyecto al cerrar esta
versión, de modo que su formulación se recoge según lo reportado por el
equipo y no según su fuente.} De él y del principio de proporcionalidad
se derivan diez requisitos:

\begin{itemize}[leftmargin=3.6em,labelwidth=3em,labelsep=0.6em,align=left,itemsep=0.2em]
  \item[M-01] Procesamiento en el borde, de modo que el fotograma no salga
        del perímetro del terminal.
  \item[M-02] No persistencia del fotograma, que se procesa en memoria y
        se descarta.
  \item[M-03] Salida agregada por zona, sin coordenadas individuales
        persistidas.
  \item[M-04] Difuminado irreversible de rostros si algún fotograma debe
        conservarse.
  \item[M-05] Plazo de retención definido, con purga automática.
  \item[M-06] Control de acceso por roles, con registro de auditoría.
  \item[M-07] Cifrado en tránsito y en reposo.
  \item[M-08] Carteles informativos conforme al Anexo~A.
  \item[M-09] Documentación del linaje de los datos.
  \item[M-10] Supervisión humana: el sistema alerta y no ejecuta acciones
        automáticas sobre personas.
\end{itemize}

Donde un parámetro admita un ajuste más protector y otro menos protector
(resolución de captura, frecuencia de muestreo, granularidad de la salida,
plazo de retención), el valor predeterminado debe ser el más protector.
Estos requisitos describen el sistema exigido y no el implementado: el
prototipo del equipo cumple M-02, pero no M-03
(Sección~\ref{sec:auditoria}).

\subsection{Requerimientos no funcionales}
\label{sec:req-no-funcionales}

Además del requisito de no identificación descrito en el
apartado~\ref{sec:privacidad-diseno}, el sistema debe cumplir cinco
requisitos no funcionales sobre su operación. La tabla recoge cada uno y
su estado verificado en el prototipo.

{\small
\setlength{\LTleft}{0pt}\setlength{\LTright}{0pt}
\begin{longtable}{@{}L{1.4cm}L{\dimexpr\linewidth-6.4cm-4\tabcolsep\relax}L{5.0cm}@{}}
\toprule
\textbf{ID} & \textbf{Requisito} & \textbf{Estado} \\
\midrule
\endfirsthead
\toprule
\textbf{ID} & \textbf{Requisito} & \textbf{Estado} \\
\midrule
\endhead
\bottomrule
\endlastfoot
RNF-02 & Servidor accesible solo desde el propio equipo, con credencial de
sesión en toda operación de escritura & Implementado \\
\addlinespace
RNF-03 & No persistir imágenes de forma permanente & Cumplido para
fuentes en vivo: los fotogramas se procesan en memoria y se descartan
(M-02). Los videos cargados para pruebas se conservan localmente en el
computador donde se ejecuta el sistema y su retención queda pendiente de
definir por el equipo \\
\addlinespace
RNF-04 & Declarar expresamente cuándo el sistema opera en modo de
demostración sintética & Implementado \\
\addlinespace
RNF-05 & No presentar ninguna cifra de desempeño sobre datos públicos ni
sobre grabaciones propias como desempeño esperado en el Aeropuerto
Internacional Jorge Chávez sin validación en sitio & Aplicado en este
informe (Secciones~\ref{sec:sintesis-antecedentes}
y~\ref{sec:analisis}) \\
\addlinespace
RNF-06 & Tiempo de inferencia compatible con la operación práctica & No
verificado como cumplido: P2PNet en CPU, sin GPU, procesó del orden de
3,5~s por cada 0,2~s de video en la prueba registrada \\
\end{longtable}
}
